\chapter{绪论}
%%====================================================================

\section{研究背景}

过去数十年间，开源软件（Open-Source Software，OSS）已从小众的技术运动成长为
现代数字基础设施不可或缺的组成部分\cite{lerner2002simple,chen2021opensource}。
Linux 操作系统支撑着全球绝大多数服务器与移动设备，
Apache HTTP Server 托管着互联网上数量庞大的网站，
Python、Git 等工具深度嵌入数百万开发者的日常工作流程——
这些软件均是开源协作模式的产物\cite{lakhani2003open}。
不同于商业软件的封闭开发，开源项目以源代码公开、全球开放协作为基本特征：
任何人均可查看代码、提交改进、报告问题，无需与项目组织存在雇佣关系。
该模式极大降低了知识传播的门槛，催生了高度去中心化的创新生态，
但也带来了一个根本性的挑战——既然贡献出于自愿，
理解驱使开发者持续投入时间和精力的因素便成为开源研究的核心议题。

上述问题在比特币社区中表现得尤为突出。
比特币是 Satoshi Nakamoto\footnote{Satoshi Nakamoto 为化名，其真实身份至今未被确认。} 于2008年提出的点对点电子现金系统\cite{nakamoto2008bitcoin,cheng2024bitcoin}，
底层协议自诞生之日起便以完全开源的形式运作，
代码库托管于 GitHub\footnote{GitHub（\url{https://github.com}）是全球最大的代码托管与开源协作平台，由微软于2018年收购。} 平台，目前已积累超过76,000名关注者，
吸引了来自全球的数千名开发者贡献代码与文档。
比特币区别于其他开源项目的关键在于：项目背后没有任何公司主体提供资金保障，
也没有雇主为核心开发者发放薪酬。项目的存续与演进完全建立在志愿贡献之上，
而这些贡献直接关乎一个数千亿美元规模金融网络的安全与可靠。
核心开发者一旦离去或社区贡献持续萎缩，比特币网络将面临严峻的技术维护风险。
因此，有效激励开发者持续贡献，是比特币社区乃至整个开源生态面临的核心治理议题。

学术界围绕开源开发者激励机制已积累了丰富成果。
早期研究发现，开发者的参与主要受内在动机驱动：对技术问题的纯粹兴趣\cite{shah2006motivation}、
在同行社区中积累声誉的渴望\cite{lerner2002simple}，以及通过贡献提升个人技能、
进而在劳动力市场获得更好机会的职业考量\cite{fang2009understanding}，
三者共同构成了驱动开发者持续投入的基础动力。
另一方面，Davidson 等人\cite{davidson2014older}指出，利他精神与自我实现感同样是重要的内在激励来源。
上述研究共同刻画出一幅图景：在开源生态中，
开发者的自主性与社区归属感往往是持续贡献的核心支撑，
外部物质回报反而居于次要地位。

不过，随着开源软件在商业和社会中的重要性持续上升，
外部经济激励机制逐渐进入研究者的视野。
货币奖励被证明能在一定程度上促进开发者贡献\cite{wang2022monetary}，
但货币激励与内在动机之间的关系颇为微妙——若管理不当，不适当的金钱激励反而可能产生“挤出效应”，
削弱开发者原有的内在驱动力\cite{benabou2003intrinsic,qiao2021mitigating}。
2019年，GitHub 正式推出 Sponsors 机制\footnote{GitHub Sponsors：\url{https://github.com/sponsors}，允许用户直接向开源开发者提供定期或一次性经济赞助。}，首次为开源贡献搭建了平台级的捐赠基础设施，
使任何人都可以向自己欣赏的开发者提供经济支持。
该机制的出现引发了学界的广泛关注。
Shimada 等人\cite{shimada2022github}发现，
部分开发者将赞助本身视为一种新型的社区贡献方式；
Zhang 等人\cite{zhang2022who}进一步揭示，赞助在短期内对开发活动具有温和的正向效应，
但获得赞助的可能性与开发者在社区中的社会地位密切相关；
Conti 等人\cite{conti2023incentivizing}通过更长期的追踪研究发现，
实际收到赞助反而可能对创新活动产生长期负向影响，
这与挤出效应的理论预测相吻合。
Overney 等人\cite{overney2020not}的研究也表明，
单纯依靠捐赠机制来提升开源贡献的效果并不显著，
暗示捐赠的效用可能远比表面上看更为复杂。

上述研究的密集涌现，折射出开源生态在现实层面所面临的深层可持续性困境。
现代数字经济的高速运转建立在大量关键开源组件之上，但这些组件的维护工作往往
高度集中于少数志愿开发者，缺乏稳定的制度性支持。
2021年底曝光的 Apache Log4j 漏洞事件\footnote{即 CVE-2021-44228，又称“Log4Shell”，被评为 CVSS 最高危险等级10.0分，影响了全球数十亿台设备。}便以极为惨烈的方式揭示了这一隐忧：
一个嵌入全球数十亿软件系统的核心日志库，长期仅由数名志愿者利用业余时间维护，
安全漏洞的存在几乎使整个互联网基础设施陷入系统性风险\cite{alami2024sustainability}。
比特币生态中，这种脆弱性以更为独特的方式呈现。
比特币网络目前承载着超过万亿美元规模的资产，
底层协议的核心维护工作却高度依赖为数不多的全职开发者；
任何核心人才的流失或社区参与的持续萎缩，
都可能对协议演进能力乃至网络整体安全产生深远影响。
在这样的背景下，以 Brink\footnote{Bitcoin Brink 官方网站：\url{https://brink.dev}。} 为代表的专项资助机构应运而生，
探索出基于信息透明化和分阶段披露的独特激励路径，
以期在不违背去中心化治理理念的前提下，维持更广泛开发者群体的参与意愿与贡献活力\cite{wu2024governance}。
对该机制进行科学评估，具有超越单一案例的现实普遍意义。

回顾这些工作可以发现，现有文献对捐赠的理解主要集中于货币转移的经济效应，
即货币激励对开发者参与的直接促进作用。
但在实际的捐赠生态中，信息的流动往往先于资金的转移，
且覆盖的受众远比实际受益者更广。
当一个资助机构在社交媒体上宣布收到一笔捐款时，
所有关注该机构的开发者都会接收到这条信息，
无论他们最终能否从中受益。这条信息本身对开发者行为产生可测量影响的可能性，
目前尚缺乏实证检验。仅凭信息而非实际资金转移所产生的激励效应，
在既有文献中几乎没有得到系统研究。
该研究空白尤其值得关注：信息披露相较于实际资金分配具有近乎为零的边际成本，
同时又能触达整个社区而非少数受资助者；
若激励效应确实存在，则将为开源社区的低成本治理提供重要的政策工具。
从信息经济学的视角审视，捐赠公告的本质是一种信号传递行为——
它向潜在贡献者同时传递了项目获得外部认可、机构具有奖励意愿以及未来奖励可得性等多重信号，
这些信号会在开发者心中激活或调整预期结构，进而影响参与决策\cite{connelly2011signaling}。
类似的逻辑在众筹研究中已有部分印证：项目进展公告能够显著提升后续的支持率，
即便公告内容并不伴随任何即时的经济转移\cite{mollick2014dynamics}。
但在开源开发者群体中，信息披露以类似机制触发可测量的行为变化，
尤其是不同类型的披露（普惠型与定向型）是否产生差异化效应，
实证层面仍是空白。

在社区内部的信息维度上，同样存在一个亟待填补的研究空白。
开源项目的 Issue 讨论区不仅是技术问题协商的场所，
也是社区成员情感状态汇聚与传导的媒介。
当项目遭遇重大 Bug、核心功能争议或对外部环境变化的集体焦虑时，
Issue 区的讨论情感会随之出现明显的极性变化。
Guzzi 等人\cite{guzzi2013}发现，开源社区的协作沟通中蕴含着丰富的情感信息，
这些信息构成了社区健康状态的重要信号。
但社区情感在时间维度上影响开发者具体行为的机制，
情感变化预测随后代码活动或讨论活动的可能性，
不同类型的开发活动之间相互促进或相互抑制的动态关系，
这些问题尚缺乏系统性的实证探讨。
现有将情感分析引入开源社区的研究，大多以描述情感分布的静态特征
或考察情感与代码质量的截面相关性为主要目标\cite{guzzi2013}，
较少采用动态因果推断框架来系统检验情感冲击在时间轴上传导至开发者的行动决策。
更为关键的是，代码活动与讨论活动之间的动态互动——即
当开发者通过编写代码回应社区问题时，是否会相应减少讨论参与，
形成活动层面的替代关系——虽有理论依据，
且对理解开源社区的协作动态具有重要意义，
却在实证研究中几乎未曾得到正面检验。
从信息行为研究的视角看，情感本身便是一种重要的信息资源，
它向社区成员传递项目状态的信号，影响行动意愿与行动方向。
循此视角理解开发者行为，有可能揭示出超越个体激励机制的社区层面动态规律。

在上述背景下，本研究将视线转向比特币开源社区，
尝试从信息视角对前述两个研究空白作出系统性的实证回应。

\section{研究场景：比特币社区与 Brink 基金会}

比特币开源社区为本研究提供了几乎难以在其他情境中复现的理想研究场景。
核心原因在于一个独特的研究对象——非盈利组织 Bitcoin Brink（以下简称“Brink”）
及其所形成的两阶段捐赠信息披露机制。

Brink 成立于2020年，是目前规模最大的专注于比特币核心协议开发的资助机构。
该机构采用众筹模式面向个人捐赠者和机构投资者募集资金，
并通过严格的申请审核流程为符合条件的比特币开发者提供全职工作资助。
基于比特币本身的去中心化哲学，Brink 刻意保持独立的非盈利定位，
资金来源和资助决策均对外公开，以维护社区信任。
正是对透明度的高度重视，催生了一套具有内在研究价值的两阶段信息披露机制。

第一阶段，Brink 收到捐款后，会在 X 平台（原 Twitter）官方账号上发布公告，
公开感谢捐赠方并注明捐款金额，但不说明资金将用于资助哪位具体的开发者。
这类公告面向所有关注者广泛传播，每一位比特币开发者都可能看到；
由于受益人尚未确定，公告传递的核心信号是"社区正在获得外部支持，
有资金可以用于奖励开发者"。本研究将这类披露定义为非定向披露（Untargeted Disclosure），
强调其普惠性——每位开发者均可合理地将自己视为潜在受益者。
第二阶段，经过通常为数周乃至数月的申请与审核流程后，
Brink 在官方网站上正式公布当期受资助开发者名单，
列出姓名、资助金额及资助期限，并对贡献做出肯定性说明。
与第一阶段不同，该类披露传递的核心信号是“谁因贡献突出而获得了认可与资金支持”。
本研究将其定义为定向披露（Targeted Disclosure），
因为它指向具体个人，对于未入选者而言也意味着明确的排除。

两次披露在时间上相互分离、在内容上存在本质差异，
这在方法论层面具有重要价值——研究者得以将两种信息效应分别识别，
而无需将信息效应与实际资金的到达混为一谈。
在本研究的观测期内（2022年1月至2023年12月），
Brink 实际资助的开发者数量极为有限，两年内仅资助了8名开发者转为全职从事比特币开发。
将这8名实际受资助者从分析样本中剔除后，
研究者便可专注考察：对于其余1610名开发者而言，
捐赠信息的公开本身（而非任何真实的金钱转移）
对活动行为产生了可测量的影响。
这种“信息与金钱的干净分离”在开源捐赠研究中极为罕见，
构成了本研究得以回答核心问题的关键条件。

捐赠信息披露机制之外，比特币社区还具备另一项研究优势：
完整、开放且可系统处理的 Issue 讨论历史。
比特币项目自创立以来积累了大量 Issue 记录，
涵盖 Bug 报告、功能讨论、安全议题乃至社区内部的技术争论——
这些文本构成了追踪社区情感动态变化的完整语料库。
结合 GH Archive 对 GitHub 活动数据的完整记录、X 平台的公告历史、
Brink 官网的资助档案以及 Google Trends 的搜索量数据，
本研究得以在同一社区中、围绕同一组开发者群体，
同时探讨外部信息与内部情感两个维度对开发者行为的影响。
数据来源的多元性与可靠性，是最终选择比特币社区作为研究对象的重要原因。

\section{研究问题}

立足上述背景与研究场景，本研究聚焦两个相互补充的核心问题。

第一个研究问题关注外部信息对开发者行为的影响。
Brink 基金会在实践中形成了非定向披露与定向披露两种信息传播方式，
二者在受众结构、信号内容和心理机制上存在本质差异，
却可能产生相互交织的激励效果。具体而言：

RQ1：探究 Bitcoin Brink 基金会的非定向捐赠披露与定向捐赠披露分别对
开发者代码活动、审查活动和讨论活动的影响，并检验当两种披露同时存在时，
二者之间是否存在交互效应。

第二个研究问题将目光转向社区内部，探讨情感信息的传导机制。
社区讨论中的情感状态不仅反映开发者的体验，
也可能作为项目状态的信号影响开发者的行动选择。
本研究检验代码活动与讨论参与这两种典型的开发活动之间是否存在情感所触发的动态路径。

RQ2：探究比特币开源社区 Issue 讨论中的情感状态在时间维度上对
开发者代码活动与讨论活动的影响，并分析代码活动的变化是否进一步影响讨论活动，
形成可识别的跨活动传导路径。

\section{研究方法概述}

针对两个研究问题，本研究采取差异化的计量策略。

研究一面对的核心方法论挑战是内生性问题：
开发者的活动水平可能反向影响外部捐款规模，
进而影响 Brink 的披露行为，在观测数据中产生反向因果偏误。
为此，本研究采用面板数据个体固定效应模型（Panel Fixed-Effects Model）
作为基准识别框架，控制个体固定效应以消除不可观测的个体异质性干扰。
在此基础上进一步引入工具变量方法（Instrumental Variable，IV），
以“Bitcoin Brink”关键词在 Google Trends 上的搜索量指数及其与定向披露的交互项
作为两个内生变量的工具变量，借助外生的公众关注度变化来识别披露金额的因果效应，
从而解决反向因果导致的估计偏差。

研究二面对的核心方法论挑战则是变量间的双向因果关系：
情感、代码活动和讨论活动三者之间可能相互影响，
单方程框架无法合理处理。
本研究采用向量自回归模型（Vector Autoregression，VAR），
将三个变量均作为内生时间序列联合建模。
在此框架下，通过格兰杰因果检验（Granger Causality Test）在统计层面识别变量间的预测关系，
并借助脉冲响应函数（Impulse Response Function，IRF）
直观呈现某一变量受到冲击后其他变量的动态响应路径，
从而系统描绘社区情感与开发活动之间的动态传导机制。

\section{研究意义}

\subsection{理论意义}

本研究最核心的理论贡献在于从信息视角重新审视开源开发者的激励机制，
将既往研究中混同于货币效应的信息效应单独提炼并加以检验。
在已有的开源激励研究中，捐赠通常被视为货币转移行为，
影响机制被归结为对开发者的直接经济补偿\cite{zhang2022who,wang2022influence}。
但在实际运作中，捐赠信息的公开传播往往先于真实的资金分配，且波及范围远大于后者，
由此形成了一种超越货币本身的信号效应。
本研究通过对样本的精心处理与工具变量的设计，
将信号效应从货币效应中剥离，首次在开源社区情境下提供了
“捐赠信息披露本身构成有效激励因素”的实证证据，
丰富了信息激励理论的层次与应用边界。

在理论框架层面，本研究将期望理论（Expectancy Theory）
与期望失验理论（Expectancy Disconfirmation Theory）联合应用于
同一研究情境的不同阶段。非定向披露通过激活“潜在受益”期望影响开发者动机，
定向披露的公布则随即在大多数非受助者中引发期望落空——
两种机制在时间上先后作用，共同塑造了开发者对信息披露的复合响应。
这种基于期望认知动态的分析框架，对于理解其他需要区分“普遍性激励”与
“选择性激励”效果的管理情境同样具有借鉴价值。

同时，关于社区情感影响路径的研究，将情感即信息理论（Affect-as-Information Theory）
引入开源开发者行为研究领域。
情感作为信息的视角强调，情感状态本身携带着关于当前处境的认知内容，
引导个体调整判断与行动方向\cite{schwarz1983moods}。
将该框架应用于开源社区，意味着把 Issue 讨论中的集体情感
理解为社区状态的公共信号，而非单纯的情绪宣泄。
另一方面，目标手段替代理论（Goal-Means Substitution Theory）被用于解释
代码活动对讨论活动的抑制效应。这一跨理论的应用拓展了两种理论在数字协作情境中的适用范围，
并为理解开源社区中不同类型活动之间的内在关联提供了新的理论工具。

从图书情报学的学科视角出发，本研究将“信息行为”这一核心议题延伸至
开源软件社区这一新兴数字协作场景，
将开发者对信息披露的响应与情感信号的加工理解为一种特殊的信息利用与决策行为。
方法论层面，本研究融合面板计量经济学与时间序列分析方法，
为图书情报学的实证研究提供了多源异构数据综合分析的示范案例，
有助于推动学科在量化研究方法上的多元化发展。

研究设计层面，本研究构建了“信息效应与货币效应分离”的识别框架，
通过将实际受资助者从样本中剔除、并以工具变量方法处理内生性问题，
在开源社区研究情境中实现了较为严格的因果推断。
这种将准自然实验思想系统融入开源行为研究的方法论路径，
有别于现有文献中普遍采用的相关性分析范式。
对于希望从“什么因素与开发者行为相关”推进到
“什么因素真正驱动了开发者行为”的后续研究者而言，
本研究提供了一个可供参照的方法论示范，
有助于推动开源社区研究从描述性分析向因果推断的范式演进。

\subsection{实践意义}

对于以 Brink 为代表的开源基金会而言，本研究揭示了信息披露方式的设计对
社区开发活力的实际影响，具有直接的政策参考价值。
非定向披露能够在整个开发者群体中激活期望，形成广泛的动机激励，
适合定期发布以维护社区的整体参与感。
但定向披露在明确受益人之后，不可避免地会在未入选者中产生失落效应，
进而削弱非定向披露的长期激励效果。
这提示基金会在信息披露时应有意识地管理披露节奏与表述方式——
例如在定向披露公告中着重强调资助计划的持续性和未来的更多机会，
以缓解未入选者的期望落空感。

对于比特币社区乃至更广泛的开源生态而言，
社区情感对开发活动的预测作用为社区健康管理提供了实践启示：
情感极性的下滑可以作为社区面临冲突或技术困境的预警信号。
在这一时期，代码活动往往随之上升，表明社区成员正在自发调动资源解决问题。
社区管理者可借助情感监测系统识别此类时间节点，
及时提供协调资源、降低开发者面临的障碍。

更广泛地看，本研究揭示的“以信息换激励”逻辑，
对所有依赖自愿贡献且资源有限的数字公共品社区均具有参考意义。
在无法依靠大规模直接资助的现实约束下，
信息的透明化与策略性披露可以作为货币激励的低成本替代或有益补充——
既维持社区信任与凝聚力，又对更广泛的潜在贡献者产生持续的参与激励。
该视角不仅适用于加密货币生态，
对于同样面临可持续性挑战的维基百科类志愿知识贡献平台、
公民科学社区以及政府推动的开放数据生态，同样具有一定的政策启发价值。

\section{论文结构安排}

